\documentclass[10pt,a4paper]{amsart}
\usepackage[margin=19mm]{geometry}
\usepackage{amsmath,amssymb,mathtools}
\usepackage[scr=rsfs,cal=euler]{mathalfa}
\usepackage{xfrac}
\usepackage[unicode,hidelinks,bookmarks=false]{hyperref}
\newcommand{\ie}{i.e.\ }
\newcommand{\cf}{cf.\ }
\newcommand{\resp}{respectively\ }
\newcommand{\Subsection}{\subsection}
\providecommand{\nobreakdash}{\nobreak\hbox{-}}
\setlength{\emergencystretch}{2em}
\multlinegap0pt
\newtheorem{proposition}{Proposition}
\def\invlim{\varprojlim}
\renewcommand{\mathcal}{\mathscr}
\newcommand{\sh}[1]{\mathcal #1}
\title{The resolution property holds away from codimension three}
\author{Siddharth Mathur, Stefan Schr\"oer}
\date{Source: arXiv:2109.09623v3 (CC BY 4.0)}
\begin{document}
\maketitle

\noindent\emph{Source and licence.} The resolution property holds away from codimension three. Version arXiv:2109.09623v3 (CC BY 4.0), distributed by arXiv under the Creative Commons Attribution 4.0 International licence, http://creativecommons.org/licenses/by/4.0/, as recorded in the arXiv licence metadata for this exact version and shown on the abstract page https://arxiv.org/abs/2109.09623v3. Full licence notice: \texttt{licenses/arxiv-2109.09623-CC-BY-4.0.txt}.\par\smallskip

\noindent\textit{Editorial scope.} Counted unit: the article's proposition proposition (source0.tex lines 438--440). The article's complete original proof of it is delivered with it (source0.tex lines 442--445, one \texttt{proof} environment of 4 lines). No mathematics is written, repaired or re-derived: the statement and the proof are the article's own lines, in the article's own notation, and no retained line is substituted or re-flowed.  No further source-local context is needed: every cross-reference of every retained line resolves inside the two delivered spans.  Nothing was omitted from any retained span, so the package carries no omission record.  No retained line contains a citation, so the wrapper carries no bibliography and no bibliography entry.  No retained line contains a figure environment, an image-inclusion command or drawing code, so no image is delivered and the package declares no asset dependency.  Editorial formatting only, in addition: an AMS \texttt{amsart} preamble that restates the article's own declarations and macros used by the retained lines, namely the environments \texttt{proposition} on separate counters, together with the standard packages those lines need.  The retained environments print in this document as Proposition 1 in order of appearance, whereas the article prints the same items with the numbering of its own full document.  The complete native source of the article as distributed in the arXiv e-print for this exact version is bundled unchanged as \texttt{sources/116244/source0.tex}.
\begin{proposition}
The natural morphism $g: U \to U'$ is an isomorphism. 
\end{proposition}

\begin{proof} There are positive integers $n$ and $m$ with $\sh{I}^n \subset \sh{I'} \subset \sh{I}$ and $\sh{I'}^m \subset \sh{I}^n \subset \sh{I}'$. Thus, if $U_n$ and $U'_m$ denote the flat neighborhoods of $Z^{(n)}$ and $Z'^{(m)}$, we obtain maps
\[U \to U' \to U_n\]
\[U' \to U_n \to U_m'\]
and if we show both compositions are isomorphisms, it will follow that $U \to U'$ is an isomorphism. Thus, it suffices to assume that $\sh{I}'=\sh{I}^n$ by symmetry. Now, the result follows from the fact that the natural map $\invlim_n (B/I^{n}) \to \invlim_m (B/I^{nm})$ is an isomorphism for any commutative ring $B$ and any ideal $I \subset B$.\end{proof}

\end{document}
